\section{Техническое задание}

\subsection{Основание для разработки}

Основанием для разработки является задание на курсовую работу по дисциплине «Проектирование информационных систем» на тему «Моделирование продукционной системы Маркова на языке программирования Python».  
Необходимость разработки данного программного продукта обусловлена потребностью в инструменте, позволяющем наглядно продемонстрировать работу формальных систем подстановок и автоматизировать процесс моделирования вычислительных алгоритмов в рамках теории формальных языков.  

В качестве нормативных документов и источников, используемых при разработке, приняты:
\begin{itemize}
	\item ГОСТ 19.102–77 «Единая система программной документации. Стадии разработки»;
	\item ГОСТ 19.201–78 «ЕСПД. Техническое задание. Требования к содержанию и оформлению»;
	\item ГОСТ 34.601–90 «Автоматизированные системы. Стадии создания»;
	\item Учебно-методические материалы кафедры по курсовому проектированию.
\end{itemize}

\subsection{Цель и назначение разработки}

Целью разработки является создание программного приложения, реализующего моделирование работы продукционной системы Маркова с возможностью визуализации процессов подстановок и анализа результатов вычислений.  

Разрабатываемое приложение предназначено для студентов, преподавателей и исследователей, изучающих принципы функционирования формальных систем, а также для демонстрации алгоритмов вычислений в учебных целях.  

Основные задачи разработки заключаются в следующем:
\begin{itemize}
	\item создание программной модели системы Маркова, включающей набор правил подстановок и входную строку;
	\item реализация интерфейса для ручного и автоматического запуска алгоритма подстановок;
	\item визуализация пошагового выполнения алгоритма, включая выделение заменяемых фрагментов строки;
	\item реализация возможности загрузки и сохранения правил подстановок из внешних файлов;
	\item ведение журнала работы (логирования) для анализа последовательности применённых правил;
	\item формирование статистики по количеству шагов, длине строки и частоте использования конкретных правил;
	\item обеспечение удобного и понятного графического интерфейса пользователя.
\end{itemize}

Результатом разработки является настольное приложение на языке программирования Python, использующее библиотеки \texttt{tkinter} и \texttt{customtkinter} для построения графического интерфейса, а также реализующее основные принципы функционирования продукционных систем.

\subsection{Требования пользователя к интерфейсу приложения}

Программное средство должно обладать интуитивно понятным интерфейсом и обеспечивать простоту взаимодействия с пользователем без необходимости предварительного изучения руководства.  

В интерфейсе должны быть реализованы следующие функции:
\begin{itemize}
	\item ввод исходной строки и правил подстановок вручную;
	\item загрузка и сохранение правил из текстовых файлов формата \texttt{.txt};
	\item запуск моделирования в двух режимах — пошаговом и автоматическом;
	\item отображение текущего состояния строки после каждого шага;
	\item подсветка применяемого правила в процессе выполнения;
	\item отображение итогового результата после завершения моделирования;
	\item сохранение истории выполнения и результатов в журнале работы;
	\item вывод статистических данных по ходу моделирования.
\end{itemize}

Архитектура приложения должна обеспечивать модульность и расширяемость. Программный комплекс состоит из следующих основных компонентов:
\begin{itemize}
	\item \texttt{MarkovSystem} — основной класс, управляющий процессом выполнения правил;
	\item \texttt{MarkovRule} — класс, описывающий отдельное правило подстановки и механизм его применения;
	\item модуль интерфейса на основе \texttt{tkinter/customtkinter};
	\item модуль логирования и статистики для анализа работы алгоритма;
	\item вспомогательные функции для загрузки, сохранения и обработки данных.
\end{itemize}

\begin{figure}[H]
	\centering
	\includegraphics[width=0.8\textwidth]{interface.png}
	\caption{Рисунок 2.1 – Композиция шаблона системы}
	\label{fig:interface}
\end{figure}

\subsection{Моделирование вариантов использования}

Для разрабатываемой программы спроектирована модель, демонстрирующая основные сценарии взаимодействия пользователя с системой. Модель отражает ключевые процессы, выполняемые приложением при работе с продукционными системами.  

Основные варианты использования включают:
\begin{enumerate}
	\item \textbf{Запуск приложения} — пользователь открывает программу, загружается главное окно.
	\item \textbf{Ввод данных} — пользователь вводит исходную строку и набор правил подстановок.
	\item \textbf{Загрузка правил из файла} — осуществляется импорт данных из текстового файла.
	\item \textbf{Пошаговое выполнение алгоритма} — система выполняет каждое правило поочерёдно с визуальным отображением изменений.
	\item \textbf{Автоматическое выполнение алгоритма} — приложение выполняет преобразование строки до достижения терминального состояния без вмешательства пользователя.
	\item \textbf{Просмотр статистики} — система выводит количество применённых правил, длину результирующей строки и другие параметры.
	\item \textbf{Сохранение отчёта} — результаты моделирования сохраняются в файл журнала.
\end{enumerate}

Для наглядности модель прецедентов представлена в виде диаграммы (рисунок~\ref{fig:usecase}), отображающей взаимодействие пользователя с основными компонентами программы.

\begin{figure}[H]
	\centering
	\includegraphics[width=0.8\textwidth]{markov_usecase}
	\caption{Диаграмма прецедентов системы моделирования}
	\label{fig:usecase}
\end{figure}

\subsection{Описание сценария использования}

Пример сценария «Пошаговое выполнение алгоритма»:
\begin{enumerate}
	\item Пользователь запускает приложение.
	\item В главном окне вводит исходную строку и правила подстановок.
	\item Нажимает кнопку «Пошаговое выполнение».
	\item Программа применяет первое правило, отображая изменение строки.
	\item Пользователь наблюдает результат и нажимает «Следующий шаг».
	\item Процесс повторяется до завершения работы алгоритма.
	\item Программа отображает итоговую строку и статистику работы.
\end{enumerate}

Предусловие: наличие корректных правил подстановок и исходной строки.  
Постусловие: программа завершает работу при достижении терминального состояния или невозможности дальнейших замен.

\subsection{Требования к оформлению документации}

Разработка программного изделия и документации должна осуществляться в соответствии с требованиями стандартов:
\begin{itemize}
	\item ГОСТ 19.102–77 «ЕСПД. Стадии разработки»;
	\item ГОСТ 34.601–90 «Автоматизированные системы. Стадии создания»;
	\item ГОСТ 19.701–90 «Схемы алгоритмов, программ, данных и систем»;
	\item ГОСТ 19.504–79 «Текст программы. Требования к содержанию и оформлению».
\end{itemize}

Документация должна включать:
\begin{itemize}
	\item титульный лист, аннотацию и оглавление;
	\item разделы с описанием архитектуры, интерфейса и сценариев работы программы;
	\item иллюстрации диаграмм классов, компонентов и вариантов использования;
	\item выводы и рекомендации по дальнейшему развитию системы.
\end{itemize}
